\noindent\textbf{DREIECKSKRANZ UM OLDENBURG.}\hfill333

\par\noindent\hspace*{\parindent}Die 21 Dreiecke und das Siebeneck haben folgende Fehler der Winkelsummen.
\begin{center}
\begin{minipage}[t]{.48\linewidth}\centering
\begin{tabular}{lrl}
Dreieck&\(21.\ 1.\ 2,\)&\(a:+0''{,}388\)\\
\guillemotright&\(21.\ 2.\ 3,\)&\(b:+5{,}973\)\\
\guillemotright&\(21.\ 3.\ 4,\)&\(c:-0{,}020\)\\
\guillemotright&\(3.\ 5.\ 4,\)&\(d:+3{,}696\)\\
\guillemotright&\(4.\ 5.\ 6,\)&\(e:+0{,}121\)\\
\guillemotright&\(5.\ 7.\ 6,\)&\(f:+1{,}142\)\\
\guillemotright&\(6.\ 7.\ 8,\)&\(g:-3{,}512\)\\
\guillemotright&\(6.\ 8.\ 9,\)&\(h:-0{,}020\)\\
\guillemotright&\(6.\ 9.\ 10,\)&\(i:+2{,}493\)\\
\guillemotright&\(9.\ 11.\ 10,\)&\(k:-0{,}798\)\\
\guillemotright&\(10.\ 11.\ 12,\)&\(l:-0{,}786\)\\
\guillemotright&\(11.\ 13.\ 12,\)&\(m:+0{,}073\)
\end{tabular}
\end{minipage}\hfill
\begin{minipage}[t]{.50\linewidth}\centering
\begin{tabular}{lrl}
Dreieck&\(12.\ 13.\ 14,\)&\(n:-1''{,}122\)\\
\guillemotright&\(12.\ 14.\ 15,\)&\(o:-0{,}866\)\\
\guillemotright&\(12.\ 15.\ 16,\)&\(p:-1{,}045\)\\
\guillemotright&\(15.\ 17.\ 16,\)&\(q:+0{,}025\)\\
\guillemotright&\(16.\ 17.\ 18,\)&\(r:-1{,}421\)\\
\guillemotright&\(16.\ 18.\ 19,\)&\(s:-4{,}920\)\\
\guillemotright&\(18.\ 20.\ 19,\)&\(t:-1{,}299\)\\
\guillemotright&\(19.\ 20.\ 21,\)&\(u:-0{,}274\)\\
\guillemotright&\(20.\ 1.\ 21,\)&\(v:-2{,}941\)\\[4pt]
\multicolumn{3}{l}{Siebeneck}\\[4pt]
\multicolumn{3}{l}{\(4.\ 6.\ 10.\ 12.\ 16.\ 19.\ 21,\quad w:+17''{,}963.\)}
\end{tabular}
\end{minipage}
\end{center}

\par\noindent\hspace*{\parindent}Hieraus ergeben sich die folgenden Bedingungsgleichungen:
\[
21\, .\,1+1\, .\,2+2\, .\,21=+0{,}388
\]
\begin{center}usw.\end{center}
\[
4\, .\,6+6\, .\,10+10\, .\,12+12\, .\,16+16\, .\,19+19\, .\,21+21\, .\,4=+17{,}963.
\]
\par\noindent\hspace*{\parindent}Hier bezeichnet \(21\, .\,1\) die an [der Richtungsbeobachtung] \(21\, .\,1\)
anzubringende Korrektur abzüglich der an \(1\, .\,21\) anzubringenden Korrektur usw.

\par\noindent\hspace*{\parindent}Bezeichnet man die Korrelaten der Bedingungsgleichungen mit den entsprechenden
Großbuchstaben \(A,B,C\) usw., so erhält man für sie die folgenden
22 Gleichungen:
\[
\begin{aligned}
-V+3A-B&=+0{,}388 & -G+3H-I&=-0{,}020\\
-A+3B-C&=+5{,}973 & -H+3I-K-W&=+2{,}493\\
-B+3C-D-W&=-0{,}020 & -I+3K-L&=-0{,}798\\
-C+3D-E&=+3{,}696 & -K+3L-M-W&=-0{,}786\\
-D+3E-F-W&=+0{,}121 & -L+3M-N&=+0{,}073\\
-E+3F-G&=+1{,}142 & -M+3N-O&=-1{,}122\\
-F+3G-H&=-3{,}512 & -N+3O-P&=-0{,}866
\end{aligned}
\]
\clearpage
\noindent334\hfill\textbf{NACHLASS.}
\[
\begin{aligned}
-O+3P-Q-W&=-1{,}045 & -T+3U-V-W&=-0{,}274\\
-P+3Q-R&=+0{,}025 & -U+3V-A&=-2{,}941\\
-Q+3R-S&=-1{,}421\\
-R+3S-T-W&=-4{,}920\\
-S+3T-U&=-1{,}299
\end{aligned}
\]
\[
-C-E-I-L-P-S-U+7W=+17{,}963,
\]
\noindent woraus sich die folgenden Werte ergeben:
\[
W=+4{,}756
\]
\[
\begin{array}{l@{\qquad}l@{\qquad}l@{\qquad}l@{\qquad}l}
A=+1{,}394&E=+3{,}305&I=+3{,}223&N=-0{,}103&R=-0{,}431\\
B=+3{,}816&F=+1{,}345&K=+1{,}489&O=+0{,}142&S=-0{,}200\\
C=+4{,}082&G=-0{,}412&L=+2{,}043&P=+1{,}394&T=-0{,}005\\
D=+3{,}694&H=+0{,}930&M=+0{,}671&Q=+0{,}329&U=+1{,}485\\
&&&&V=-0{,}021.
\end{array}
\]
\par\noindent\hspace*{\parindent}Die Korrekturdifferenzen hängen damit folgendermaßen zusammen:

\par\noindent\hspace*{\parindent}Für eine Außenseite gilt \(1\, .\,2=A\). Für eine Innenseite gilt \(21\, .\,4=W-C\). Für
eine Zwischenseite gilt \(21\, .\,1=A-V\). [Bezeichnet \((1\, .\,2)\) die Korrektur der
Richtung von 1 nach 2 usw., so gilt \(+(1\, .\,2)=-(2\, .\,1)=\frac12\times1\, .\,2=\frac12A\),
\((21\, .\,4)=-(4\, .\,21)=\frac12W-\frac12C\) usw.]

\par\noindent\hspace*{\parindent}Diese Werte wurden übrigens bereits zugrunde gelegt, um die vorstehende
Tabelle [S.\ 331/332, 1.\ Spalte], soweit dies mit 3 Dezimalstellen möglich ist,
absolut zu orientieren.

\par\noindent\hspace*{\parindent}Direkt angewandt ergeben diese Korrelaten eine neue, einmal ausgeglichene Tabelle,
[die in der ersten Zusammenstellung in der Spalte \guillemotright1. Ausgleichung\guillemotleft{}
enthalten ist. Aus ihr ergibt sich die folgende Dreiecksübersicht:]
